\section{Before Starting the Campaign}

Here are some tips and suggestions intended for the GM. Following them is not mandatory, but it might help you run the game.

\subsection{Required Knowledge}

Since the story does not touch on any events from \bookSA{} and the characters start from Shinovar, they are not required to know much about Roshar world or story presented in the books.

Of course, you as a Game Master will need to know rules of \cosmereRPG{} and enough about Roshar world to paint the picture for your players.

\subsection{Read the Adventure}

Unless you mastered improvisation, spend some time to familiarize yourself with the adventure. It will help you adjust the story to players' characters and their personal stories, as well as it will allow you to quietly foreshadow key elements of the plot.

\subsection{Players’ Expectations}

Inform players that the adventure will revolve around Reshi Isles and that Radiant powers will be limited. Discuss with players what are their expectations for the game. Consider following questions:
\begin{itemize}
	\item How serious should the adventure be?
	\item What topics players would prefer to avoid during the game?
	\item Is there anything the players would like to see in the game?
	\item Are the players comfortable with their characters dying?
	\item Are the players comfortable with characters' romance, nudity or other sexual content?
	\item If you are of age - should your group allow alcoholic drinks?
	\item What rule variants should be in effect?
	\item How often should your group meet and how long should a single game session take?
	\item How should you handle situations when one player cannot attend a game session?
\end{itemize}

Gathering answers to those (and similar) questions is not a requirement, but it will tell you about players’ hopes for the game and it will help you create a better experience for the whole group.

\subsection{Create Player Characters}

Unless your players have their characters ready after previous adventures, they must create a new party. This book contains several pre-made character sheets you can use in your game (see \hyperref[chap:exampleSheets]{Appendix \ref{chap:exampleSheets}}.), but you should encourage players to create more personalized heroes for the story.

While creating new characters, you might want to follow those guidelines:
\begin{itemize}
	\item encourage Human Ancestry, unless the player is familiar more complex Singer rules and culture.
	\item encourage Shin culture, unless the player is already familiar with the world of Roshar.
	\item another cultures relevant for the story include: Iriali, Riran, Reshi, Herdazian and Thaylean.
	\item share \hyperref[backstoryIdeas]{\textbf{Character Backstory Ideas table}} with your players for possible inspirations for characters, their motivations and goals for the story.
	\item warn players that Radiant Paths might be more challenging to play this adventure. While the story does not prevent playing with those, some of their powers are limited and players should be aware of that before choosing those paths.
	\item if a player plans playing as Knight Radiant, make sure you know what spren to introduce in the story.
\end{itemize}

\subsection{Character Relationships}

Allow players do determine how their characters know each other and what is the reason they are traveling together. Establishing characters' relations before the adventure start will make it less likely the heroes will be motivated to split the party and follow their own goals.

You can add more details to character's relationship by requesting each player to ask about some specific details two other characters. Players can think of their own questions they want to ask each other or they can choose something from the list of below:

\begin{itemize}
	\item "Once you shared with me your strange, irrational fear you never told anyone else. What was it?"
	\item "We have one thing in common what we would never guess we could share. What is it?"
	\item "Once you told me you wanted to travel with me to some remote location. Where would that be?"
	\item "You often call me a silly name. What's the name and why did you start using it?"
	\item "Your religious practices differ from those commonly observed in your lands. Why?"
	\item "I noticed you play with some small item. What is it and why do you carry it?"
	\item "You told me once that you lost someone close to you. Who was it and what would you give to get them back?"
	\item "You once told me there are worse fates than death. What did you have in mind?"
	\item "What is the case that you are willing to die for?"
\end{itemize}

\subsection{Character Backstory}

Encourage players to write a short character backstory to better root them in the world and the story. Consider if any NPCs from those backstories could make an appearance in this adventure and how they can enrich the story. You can modify scenes described in this book to include those NPCs, for example:
\begin{itemize}
	\item additional NPC can join the expedition and be lost on the island
	\item an NPC known to players’ characters might be sick and tended by Melio
	\item Reshi village might be joined by someone from the player’s backstory
	\item an old friend might be met while searching for aid
\end{itemize}


